In this work, we consider an adversarial setting modeled by a simplified version of the Dolev-Yao attacker model~\cite{Dolev-Yao}. The attacker has the capability to inject, drop, or modify packets transmitted over the network. Specifically, the attacker can introduce forged packets into the communication channel with the intent of deceiving the receiver or disrupting the communication protocol. The attacker can also intercept and remove packets from the communication channel, preventing their delivery to the intended recipient. Additionally, the attacker can alter the content of transmitted packets, potentially corrupting the data or causing the receiver to accept false information.

We assume that the attacker does not possess the cryptographic keys used for securing the communication. Consequently, the attacker cannot decrypt encrypted messages or generate valid encrypted messages and authentication tags. The inability to decrypt or correctly encrypt messages limits the attacker to manipulate the transmission medium and observable packet structures without understanding or creating valid encrypted content.

This threat model focuses on preserving the integrity and authenticity of messages in the presence of an active adversary capable of tampering with the communication channel but lacking the cryptographic keys necessary to compromise confidentiality or generate valid cryptographic tokens. 